\documentclass[11pt]{article}
\usepackage[T1]{fontenc}
\usepackage{lmodern}
\usepackage[margin=1in]{geometry}
\usepackage{amsmath,amssymb,amsthm,mathtools}
\usepackage{microtype}
\usepackage[colorlinks=true,allcolors=blue]{hyperref}
\hypersetup{pdftitle={The Sperner property of Artinian complete intersections in characteristic zero},pdfauthor={Alec Kriebel},pdfsubject={An immediate consequence of Artinian EGH}}
\setlength{\emergencystretch}{2em}
\newtheorem{theorem}{Theorem}
\newtheorem{lemma}[theorem]{Lemma}
\newtheorem{proposition}[theorem]{Proposition}
\DeclareMathOperator{\Ann}{Ann}
\DeclareMathOperator{\gr}{gr}
\newcommand{\m}{\mathfrak m}
\newcommand{\N}{\mathbb Z_{\ge0}}
\title{The Sperner property of Artinian complete intersections\\in characteristic zero}
\author{Alec Kriebel\\\small\href{https://orcid.org/0009-0001-9320-500X}{ORCID: 0009-0001-9320-500X}}
\date{October 6, 2026}
\begin{document}
\maketitle
\begin{abstract}
We record the Sperner property for every standard graded Artinian complete
intersection over an arbitrary characteristic-zero field: the maximum
minimal number of generators of an ideal, including nonhomogeneous ideals,
equals the largest coefficient of its complete-intersection Hilbert
series. The result is an immediate consequence of OpenAI's Artinian
Eisenbud--Green--Harris theorem and the previously published implication
of Harima, Wachi, and Watanabe. We give an explicit account of that
implication, eliminate degree-one generators, justify the passage to all
ideals by the associated graded filtration, and verify the field scope.
No new proof of EGH or new Lefschetz assertion is claimed.
\end{abstract}

\section{Statement and attribution}
Let \(k\) be a field and let \(A\) be a standard graded Artinian
\(k\)-algebra with \(A_0=k\). Its unique maximal ideal is
\(\m=A_{>0}\), and Nakayama's lemma gives
\(\mu_A(I)=\dim_k(I/\m I)\) for every ideal \(I\).
The Dilworth number is \(\max_{I\triangleleft A}\mu_A(I)\).
The algebra has the \emph{Sperner property} when this equals
\(\max_j\dim_k A_j\). The ideals in this definition include \(0\),
\(A\), and nonhomogeneous ideals.

\begin{theorem}\label{thm:main}
Let \(k\) be any characteristic-zero field, let \(n\ge0\), and let
\[
 A=k[x_1,\ldots,x_n]/(f_1,\ldots,f_n),
 \qquad d_i=\deg f_i\ge1,
\]
where the variables have degree one and the \(f_i\) form a homogeneous
regular sequence. Then
\begin{equation}\label{eq:target}
 \max_{I\triangleleft A}\mu_A(I)
 =\max_{j\ge0}\dim_k A_j
 =\max_{j\ge0}[t^j]\prod_{i=1}^n(1+t+\cdots+t^{d_i-1}).
\end{equation}
If \(p\) is the last degree attaining the maximum, equality is attained
by \(\m^p\). For \(n=0\), or when all \(d_i=1\), the algebra is \(k\),
the product is \(1\), and the maximum is \(1\), attained by \(A\).
\end{theorem}

Harima--Wachi--Watanabe \cite[Theorem~11]{HWW} proved that EGH for a
given graded complete intersection implies its Sperner property. Their
Definition~2 already ranges over all ideals; their Proposition~8 invokes
Watanabe's graded-ideal reduction \cite[Lemma~2.4]{Watanabe}.
OpenAI's characteristic-zero EGH theorem \cite[Corollary~1.3]{OAI:EGH}
supplies the previously conditional input. Its companion
\cite[Theorem~1.1 and Corollary~1.2]{OAI:LPP} proves stronger Betti
inequalities, which are not needed here.
This note records the resulting consequence and checks its scope;
neither the upstream breakthrough nor the known implication and
reductions are claimed as new work of this author.
The upstream manuscripts are dated September~23, 2026; the public
collection announcement is dated October~6, 2026 \cite{Release}.
No claim of first priority is made.

\section{The exact EGH input and scope reductions}
The only EGH assertion needed is the following Hilbert-function statement.
For \(S=k[z_1,\ldots,z_q]\) and a full homogeneous regular sequence
\(g_1,\ldots,g_q\) of degrees \(a_i\ge2\), every homogeneous ideal
\(Q\supseteq(g_1,\ldots,g_q)\) has a homogeneous counterpart
\(Q'\supseteq(z_1^{a_1},\ldots,z_q^{a_q})\) such that
\begin{equation}\label{eq:EGH}
\dim_k(S/Q)_j=\dim_k(S/Q')_j\qquad\text{for every }j\ge0.
\end{equation}
We may sort the degrees, since every permutation of a homogeneous system
of parameters in the Cohen--Macaulay ring \(S\) is regular; relabeling
the pure-power variables then restores the original indexing.
The counterpart may have extra generators in degrees beyond those of
\(Q\); no assertion about its generator count is required.
OpenAI \cite[Theorem~1.1 and Corollary~1.2]{OAI:EGH} gives this over
\(\mathbb C\) using commuting linear forms with division-algebra
coefficients, an ordered basis, and monomial extraction. We invoke that
external theorem explicitly, rather than deriving it in this note.

For completeness, its full-length field descent needs no embedding of
an arbitrary \(k\) into \(\mathbb C\). For a fixed \(Q\), choose a
finitely generated coefficient field \(k_0\subseteq k\) over
\(\mathbb Q\) containing coefficients of generators of \(Q\) and of
the \(g_i\). Include the \(g_i\) in the chosen generators over \(k_0\).
Regularity descends to \(k_0\) by faithful flatness. Embed \(k_0\) into
\(\mathbb C\) and apply the complex theorem after base change.
The resulting counterpart can be chosen monomial. Use the same monomial
exponents to define it over \(k_0\) and then over \(k\).
Dimensions of graded pieces are preserved by both field extensions from
\(k_0\), proving \eqref{eq:EGH} over \(k\).
Only full-length EGH is used; no general-length Artinian reduction is
needed.

If the original sequence has \(r\) linear members, these are linearly
independent: dependence would give fewer than \(n\) generators for an
ideal of height \(n\), contradicting the height theorem. A graded
linear coordinate change eliminates them and gives
\[
 A\simeq k[z_1,\ldots,z_{n-r}]/(g_1,\ldots,g_{n-r}).
\]
The remaining forms retain their degrees, all at least two. Their
Artinian quotient makes them a homogeneous system of parameters in a
polynomial ring; Cohen--Macaulayness makes them a regular sequence.
The graded isomorphism preserves all ideals, the numbers \(\mu_A(I)\),
and all Hilbert dimensions. Each eliminated factor in
\eqref{eq:target} is \(1\). If \(r=n\), the case \(A=k\) is immediate.

Successive regular-sequence exact sequences give
\begin{equation}\label{eq:hilbert}
 \sum_j h_jt^j
 =\prod_i(1+t+\cdots+t^{d_i-1}),
 \qquad h_j=\dim_k A_j.
\end{equation}
Put \(c=\sum_i(d_i-1)\). Artinian complete-intersection duality gives
\(A_c\simeq k\) and perfect multiplication pairings
\(A_j\times A_{c-j}\to A_c\); equivalently the Koszul resolution
gives the graded canonical module \(A(c)\).
The Hilbert function is symmetric and unimodal. An elementary explanation
of unimodality is the symmetric chain decomposition used next.

\section{Matching and the homogeneous bound}
We reproduce the mechanism of \cite[Propositions~7--8 and Theorem~11]{HWW}.
Let
\[
 B=k[z_1,\ldots,z_q]/(z_1^{a_1},\ldots,z_q^{a_q}).
\]
Its monomial basis is the ranked product of chains
\(0\le e_i\le a_i-1\).
Such a product admits a symmetric saturated chain decomposition.
One direct induction starts with the rectangle \(0\le i\le a\),
\(0\le j\le b\), where \(a\le b\), and uses the chains
\[
 (r,0),(r,1),\ldots,(r,b-r),(r+1,b-r),\ldots,(a,b-r),
 \qquad 0\le r\le a.
\]
They partition the rectangle and their endpoint ranks sum to \(a+b\).
Applying this construction to each old chain times a new factor preserves
symmetry about half the total rank. Rank sizes are therefore unimodal.
Whenever \(h_j\le h_{j+1}\), each chain meeting rank \(j\) extends
to rank \(j+1\): below the middle this follows from symmetry; at or
above the middle no new chain starts, so an ending would force a strict
decrease. Chain successors inject any rank-\(j\) subset into its
upper shadow.

Consequently, for every vector subspace \(W\subseteq B_j\),
\begin{equation}\label{eq:matchingB}
 h_j\le h_{j+1}\quad\Longrightarrow\quad
 \dim_k W\le\dim_k B_1W.
\end{equation}
Indeed, echelonize a basis of \(W\) by a multiplicative monomial order.
Its leading monomials form a set of size \(\dim W\). Each surviving
variable multiple of a leading monomial is the leading monomial of the
corresponding product. Its upper shadow thus lies in the leading
monomial set of \(B_1W\), and the chain injection proves the inequality.

For \(V\subseteq A_j\), apply \eqref{eq:EGH} to the preimage of
the ideal \(AV\) in the polynomial ring, obtaining an ideal \(J\) of
\(B\) with \(H_{A/AV}=H_{B/J}\). The algebras \(A,B\) have the same
Hilbert function, and standard grading gives
\[
 \dim J_j=\dim V,\qquad \dim J_{j+1}=\dim A_1V.
\]
For \(V=0\) the conclusion is immediate. Otherwise, if
\(h_j\le h_{j+1}\), \eqref{eq:matchingB} and the ideal property give
\begin{equation}\label{eq:matchingA}
 \dim V=\dim J_j\le\dim B_1J_j\le\dim J_{j+1}=\dim A_1V.
\end{equation}
This is the matching property of \(A\). In particular, Hilbert matching
has not been used to assert equality of ideal-generator counts.

Let \(p\) be the last maximal Hilbert degree. Matching holds in every
degree \(j<p\), and symmetry gives \(c-p\le p\). If a homogeneous
ideal \(I\) has initial degree \(\alpha<p\), its truncation
\(I'=\bigoplus_{j\ge\alpha+1}I_j\) is an ideal and
\begin{equation}\label{eq:truncate}
 \mu(I')-\mu(I)=\dim A_1I_\alpha-\dim I_\alpha\ge0.
\end{equation}
This follows by counting the degree pieces of \(I/\m I\).
Iterating, including the zero-ideal case, yields
\begin{equation}\label{eq:compression}
 \mu(I)\le\mu(I\cap\m^p).
\end{equation}

Now assume \(I\subseteq\m^p\), and set \(J=\Ann_A(\m I)\).
Gorenstein duality gives
\(\dim\Ann(U)=\dim A-\dim U\) for every ideal \(U\).
To see the equality with the orthogonal complement, compose
multiplication with projection onto \(A_c\). A vector orthogonal to
the ideal \(U\) is orthogonal to every multiple of its elements;
nondegeneracy then shows that it annihilates them.
Since \(\m J\subseteq\Ann(I)\),
\begin{equation}\label{eq:annmu}
 \mu(J)\ge\dim\Ann(\m I)-\dim\Ann(I)=\mu(I).
\end{equation}
The graded perfect pairings also give
\(\Ann(\m^u)=\m^{c-u+1}\) for \(0\le u\le c+1\):
one inclusion follows by degree, and any nonzero lower-degree element
pairs nontrivially with an element of \(\m^u\).
Therefore
\[
 J\supseteq\Ann(\m^{p+1})=\m^{c-p}\supseteq\m^p.
\]
Applying \eqref{eq:compression} to \(J\) gives
\(\mu(I)\le\mu(J)\le\mu(\m^p)=h_p\).
Applying it first to any homogeneous \(I\) proves the homogeneous bound.

\section{All ideals and equality}
For an arbitrary ideal \(I\), use the finite filtration
\(F^jA=\m^j=\bigoplus_{u\ge j}A_u\), and intersection filtrations
on \(I\) and \(\m I\). Under the natural identification
\(\gr_\m A\simeq A\), the initial ideal
\[
 J=\bigoplus_j (I\cap\m^j)/(I\cap\m^{j+1})
\]
is homogeneous. Its pieces are the lowest-degree initial parts of
elements of \(I\). Multiplication of such initial parts shows
\(\m J\subseteq\gr_F(\m I)\). Finite associated grading preserves
total dimensions, so
\[
 \mu(I)=\dim J-\dim\gr_F(\m I)
 \le\dim J-\dim\m J=\mu(J)\le h_p.
\]
Equality \(\gr_F(\m I)=\m J\) is neither assumed nor necessary.
Finally \(\m^p/\m^{p+1}\simeq A_p\), so \(\m^p\) attains \(h_p\).
This proves Theorem~\ref{thm:main}.

The generator bound itself also descends along a field extension
\(K/k\), if one is used: flatness gives
\[
 (IA_K)/(\m_KIA_K)\simeq (I/\m I)\otimes_kK,
 \qquad \mu_{A_K}(IA_K)=\mu_A(I),
\]
and \(\dim_K(A_K)_j=h_j\). An upper bound over \(K\) descends to
every ideal over \(k\); the equality witness already exists over \(k\).
No descent of arbitrary ideals from \(K\) is asserted.

\section{Examples and limitations}
The following Hilbert functions are coefficient lists in increasing degree.
They illustrate \eqref{eq:target}, rather than serve as evidence for EGH.
\begin{center}
\begin{tabular}{c|l|c}
degrees & Hilbert function & maximum\\\hline
\(\varnothing\) or \((1,1,1)\) & \(1\) & \(1\)\\
\((1,2,4)\) & \(1,2,2,2,1\) & \(2\)\\
\((2,2,2)\) & \(1,3,3,1\) & \(3\)\\
\((2,3,4)\) & \(1,3,5,6,5,3,1\) & \(6\)\\
\((3,3,3)\) & \(1,3,6,7,6,3,1\) & \(7\)\\
eight quadrics & \(1,8,28,56,70,56,28,8,1\) & \(70\)
\end{tabular}
\end{center}
For \(n\) quadrics the maximum is \(\binom n{\lfloor n/2\rfloor}\).
For instance, in the eight-quadric case every ideal has at most
\(70\) generators and \(\m^4\) has exactly \(70\).
The accompanying integer program reproduces the examples by convolution
and independent enumeration of bounded exponent vectors.

The theorem concerns characteristic zero and standard graded Artinian
complete intersections. It supplies no weak or strong Lefschetz property,
no assertion for nongraded complete intersections, and no unrestricted
positive-characteristic assertion. The downstream conditional argument
and monomial matching work over any field; the restriction comes from
the general EGH input used here.

\paragraph{Disclosure and verification.}
AI tools were used extensively in research, drafting, computation and
verification. The mathematical input from OpenAI is itself attributed
to OpenAI using the supplied manuscript citations. Internal automated
adversarial audits are documented in the supplement; neither they nor
the calculations constitute conventional human peer review. This
preprint has not undergone conventional human peer review or refereeing
at publication. No formalization of this follow-on theorem is claimed.
The source/dependency hashes and reproduction instructions accompany it.

\begin{thebibliography}{9}
\bibitem{HWW}
T. Harima, A. Wachi and J. Watanabe,
\emph{The EGH Conjecture and the Sperner property of complete intersections},
Proc. Amer. Math. Soc. \textbf{145} (2017), no.~4, 1497--1503.
\href{https://doi.org/10.1090/proc/13347}{doi:10.1090/proc/13347};
\href{https://arxiv.org/abs/1601.06928v1}{arXiv:1601.06928v1} (2016).

\bibitem{Watanabe}
J. Watanabe, \emph{The Dilworth Number of Artinian Rings and Finite Posets
with Rank Function}, Adv. Stud. Pure Math. \textbf{11} (1987), 303--312.
\href{https://doi.org/10.2969/aspm/01110303}{doi:10.2969/aspm/01110303}.

\bibitem{OAI:EGH}
OpenAI, \emph{Commuting Division-Coefficient Forms and the Artinian
Eisenbud--Green--Harris Conjecture}, OpenAI Math Release preprint,
September~23, 2026.
\href{https://github.com/openai/math/blob/main/preprints/Commuting-Division-Coefficient-Forms-and-the-Artinian-Eisenbud-Green-Harris-Conjecture-September-23-2026/paper.pdf}{OAI:Commuting-Division-Coefficient-Forms-and-the-Artinian-Eisenbud-Green-Harris-Conjecture-September-23-2026}.
Audited version: repository commit
\texttt{adc7f1241b42e322a6451854ab7e4b4c146bf78a}.

\bibitem{OAI:LPP}
OpenAI, \emph{The Artinian Lex-Plus-Powers Betti Theorem},
OpenAI Math Release preprint, September~23, 2026.
\href{https://github.com/openai/math/blob/main/preprints/The-Artinian-Lex-Plus-Powers-Betti-Theorem-September-23-2026/paper.pdf}{OAI:The-Artinian-Lex-Plus-Powers-Betti-Theorem-September-23-2026}.
Audited version: repository commit
\texttt{adc7f1241b42e322a6451854ab7e4b4c146bf78a}.

\bibitem{Release}
OpenAI, \emph{Sharing AI progress in mathematics}, October~6, 2026.
\url{https://openai.com/index/sharing-ai-progress-in-mathematics/}.
\end{thebibliography}
\end{document}
